% READ-TWICE-Q4-6
\section*{READ-TWICE-Q4-6: rereading collapses the Q4 edge-cover DAG}

\textbf{STATUS: PROVED construction plus exact finite verification; global minimum size remains OPEN. P versus NP remains OPEN.}

For full genuine 3-SAT on four variables, the 32 clause-presence bits identify with the
32 edges of $Q_4$. A selected clause set is UNSAT exactly when its falsifying edges cover
all 16 assignments. Hence
\[
\mathrm{UNSAT}(X)
=
\bigwedge_{v\in V(Q_4)}
\bigvee_{e\ni v} x_e.
\]
This is a monotone CNF with 16 width-four clauses, and every edge variable occurs exactly
twice.

\paragraph{Explicit read-twice BP.}
Evaluate the 16 vertex clauses sequentially. At vertex $v$, query its four incident edge
bits until one is 1; reject if all four are 0. This uses
\[
16\cdot4=\boxed{64}
\]
query nodes. Each edge variable occurs once in each endpoint gadget, hence at most twice on
every root-sink path. The program is therefore syntactic read-twice.

This directly exhibits the splitting power emphasized in the bounded-repetition branching-program
literature: one extra permitted read can turn separately simple constraint blocks into a very
small global program.

\paragraph{Exact weight-eight verification.}
All
\[
\binom{32}{8}=10{,}518{,}300
\]
weight-eight inputs were enumerated. The program has zero mismatches with direct edge-cover
testing. Exactly 272 weight-eight inputs are edge covers; independent recursion confirms that
these are exactly the 272 perfect matchings of $Q_4$.

After optimizing the local incident-edge orders, the tested program has:
\[
\text{worst promised query count}=41,
\]
\[
\text{worst accepting query count}=40,
\qquad
\text{worst rejecting query count}=41.
\]
All 272 perfect matchings take exactly 40 queries.

\paragraph{Architecture-specific lower bound.}
For any fixed local ordering of the four incident edges at each vertex, symmetry of perfect
matchings makes the selected matching edge uniform among the four incident edges at each
vertex. Thus its expected local query position is
\[
\frac{1+2+3+4}{4}=\frac52.
\]
Summing over 16 vertices gives average perfect-matching cost
\[
16\cdot\frac52=40.
\]
Therefore every vertex-by-vertex local-clause architecture has some accepting path of length
at least 40. The tested ordering attains 40 on every perfect matching and is within one query
of this architecture-specific lower bound over the entire weight-eight promise.

\paragraph{Collapse from read-once to read-twice.}
The best currently constructed adaptive FBDD for the same full clause-presence function has
2432 nodes. The explicit read-twice BP gives
\[
2432\longrightarrow\boxed{64}.
\]
The 2432 number is only an FBDD upper bound, so this is a separation between current
constructions, not a proved lower bound gap. Nevertheless the 64-node construction is exact
and demonstrates that rereading is a major escape mechanism for this laboratory.

\paragraph{Why this does not generalize directly.}
The read-twice property is special to $n=4$. A genuine 3-clause on $n$ variables is falsified
by $2^{n-3}$ assignments. In the analogous assignment-cover CNF, each clause-presence variable
therefore occurs $2^{n-3}$ times. Clause-by-clause evaluation becomes syntactic
read-$2^{n-3}$ rather than read-twice. Thus this construction is a finite hostile control,
not an asymptotic algorithm for general 3-SAT.

\paragraph{Next target.}
Attempt a hybrid STORE/RECOMP construction below 64 nodes. Equivalently, compute a
$(1,+k)$ curve in which only $k$ edge variables may be reread and the remaining endpoint
information must be retained in the DAG state. Measure node count, total path queries,
distinct-bit queries, and retained-state width separately.

\begin{thebibliography}{9}
\bibitem{BRS93}
A.~Borodin, A.~Razborov, and R.~Smolensky,
\emph{On Lower Bounds for Read-$k$-Times Branching Programs},
Computational Complexity 3:1--18, 1993.

\bibitem{Jukna95}
S.~Jukna,
\emph{A Note on Read-$k$ Times Branching Programs},
RAIRO Theoretical Informatics and Applications 29(1):75--83, 1995.

\bibitem{SavickyZak00}
P.~Savicky and S.~Zak,
\emph{A Read-Once Lower Bound and a $(1,+k)$-Hierarchy for Branching Programs},
Theoretical Computer Science 238:347--362, 2000.
\end{thebibliography}

% END_OF_NOTE::READ_TWICE_Q4_6
